\section{Conclusion}
This study demonstrates that large language models can effectively personalize POI descriptions to enhance user engagement and perceived information quality. Through a controlled user study with 51 participants evaluating 510 paired comparisons across 10 diverse POIs, we found that AI-generated adaptations significantly improved perceived clarity ($p = 0.0218$) and showed marginal improvements in trustworthiness ($p = 0.050$) compared to original manual descriptions. AI-personalized descriptions shifted user responses from neutral assessments toward stronger positive endorsements across all evaluated dimensions, while maintaining engagement levels comparable to carefully crafted manual content.

The findings reveal three key insights. First, personalization helps users form more confident and favorable assessments of POI value by surfacing information aligned with their interests and travel styles. Second, thoughtfully executed AI personalization can enhance rather than undermine perceived trustworthiness, challenging concerns about credibility of generated content. Third, participants expressed high comfort with AI-generated travel content (mean = 3.94 on a 5-point scale), suggesting receptiveness to this technology in tourism contexts.

These results have practical implications for travel platforms seeking to scale personalized experiences. LLM-based description adaptation offers a path to deliver customized content without the resource constraints of manual authoring for diverse user segments. However, deployment should consider demographic differences in receptiveness to AI content and implement strategies to prevent over-personalization from limiting serendipitous discovery.

This research contributes to the growing understanding of AI-powered personalization in tourism by demonstrating that content adaptation, not just recommendation, plays a valuable role in enhancing user experience. As large language models continue to advance, their application to travel content personalization presents opportunities to make tourism information more accessible, relevant, and engaging for diverse global audiences.
